%!TEX TS-program = xelatex

\documentclass[aspectratio=169]{beamer}

\newbool{russian}
\booltrue{russian}  % Подключаем русскоязычный вариант логотипа и надписей

% Подключаем стиль презентации HSE:
\usepackage{HSE-theme/beamerthemeHSE}

% Работа с русским языком и шрифтами
\usepackage[english,russian]{babel}
\usepackage[no-math]{fontspec}
\setsansfont{HSESans-Regular}[
    Path = HSE_Sans/,
    BoldFont = HSESans-Bold.otf,
    ItalicFont = HSESans-Italic.otf
]
\setmonofont{Courier New}
\usepackage{mathspec}
\setmathsfont(Digits,Latin,Greek)[
  Path = HSE_Sans/,
  Numbers = {Lining,Proportional}
]{HSESans-Regular}
\setmathrm[Path = HSE_Sans/,Numbers={Lining,Proportional}]{HSESans-Regular}

\uselanguage{russian}
\languagepath{russian}
\deftranslation[to=russian]{Theorem}{Теорема}
\deftranslation[to=russian]{Definition}{Определение}
\deftranslation[to=russian]{Definitions}{Определен��я}
\deftranslation[to=russian]{Corollary}{Следствие}
\deftranslation[to=russian]{Fact}{Факт}
\deftranslation[to=russian]{Example}{Пример}
\deftranslation[to=russian]{Examples}{Примеры}

\usepackage{graphicx}   % Для вставки изображений
\graphicspath{{images/}} % Папка с картинками (если будете использовать)

%%%%%%%%%%%%%%%%%%%%%%%%%%%%%%%%%%%%%%%%%%%%%%%%%%%%%%%%%%%%
% Информация о выступлении

\title[Библиотека на C++]{%
  \fontsize{12}{14}\selectfont Разработка масштабируемой библиотеки для анализа и обработки графов с\\
  \fontsize{12}{14}\selectfont поддержкой многопоточных вычислений на языке C++\\[0.3cm]
  \fontsize{10}{12}\selectfont Development of a scalable library for graph analysis and processing with\\
  \fontsize{10}{12}\selectfont support for multithreaded computing in C++
}

\subtitle{\fontsize{9}{11}\selectfont Индивидуальный программный проект}

\author[Романюков Е.С.]{%
  2 курс ОП «Прикладная математика и информатика», группа БПМИ237\\
  \textbf{Романюков Егор Сергеевич}\\[0.3cm]
  Научный руководитель:\\
  \textbf{Галицкий Борис Васильевич}, Приглашенный преподаватель ФКН НИУ ВШЭ
}

\institute[НИУ ВШЭ]{%
  Национальный исследовательский университет\\
  «Высшая школа экономики»
}

\date{}
%%%%%%%%%%%%%%%%%%%%%%%%%%%%%%%%%%%%%%%%%%%%%%%%%%%%%%%%%%%%

\begin{document}

% Титульный слайд
\frame[plain]{\titlepage}

%%%%%%%%%%%%%%%%%%%%%%%%%%%%%%%%%%%%%%%%%%%%%%%%%%%%%%%%%%%%
% Разделы и слайды по структуре презентации

% 1. Описание предметной области
\section{Описание предметной области}
\begin{frame}
  \frametitle{Описание предметной области}
  \begin{itemize}
    \item \textbf{Графы} — фундаментальные структуры данных, представляющие наборы объектов и связи между ними
    \item \textbf{Применение графов}
      \begin{itemize}
        \item Социальные сети (связи между пользователями)
        \item Дорожные сети (оптимальные маршруты)
        \item Web-графы (структура веб-страниц)
        \item Биологические сети (взаимодействие молекул)
        \item Телекоммуникации (маршрутизация данных)
      \end{itemize}
    \item \textbf{Алгоритмы на графах} решают ключевые задачи в науках о данных и программной инженерии
  \end{itemize}
\end{frame}

% 2. Актуальность работы
\section{Актуальность работы}
\begin{frame}
  \frametitle{Актуальность работы}
  \begin{itemize}
    \item \textbf{Рост объема данных} требует эффективных инструментов для их анализа
    \item \textbf{Современные задачи} требуют работы с большими графами (миллионы вершин)
    \item \textbf{Существующие решения} часто
      \begin{itemize}
        \item Имеют ограниченную поддержку многопоточности
        \item Не полностью используют возможности современного C++
        \item Требуют дополнительных зависимостей от C-библиотек
        \item Имеют ограниченный функционал
      \end{itemize}
    \item \textbf{Потребность} в производительной, гибкой и удобной библиотеке для работы с графами
  \end{itemize}
\end{frame}

% 3. Цель и задачи
\section{Цель и задачи}
\begin{frame}
  \frametitle{Цель и задачи}
  \textbf{Цель}
  \begin{itemize}
    \item Разработать универсальную библиотеку для работы с графами на C++ с поддержкой многопоточных вычислений
  \end{itemize}

  \vspace{0.5cm}
  \textbf{Задачи}
  \begin{enumerate}
    \item Реализовать различные способы представления графов
    \item Имплементировать набор классических алгоритмов на графах
    \item Разработать эффективные параллельные алгоритмы
    \item Провести комплексное тестирование библиотеки
    \item Обеспечить подробную и понятную документацию
    \item Создать инструменты для визуализации графов
  \end{enumerate}
\end{frame}

% 4. Анализ существующих решений
\section{Анализ существующих решений}
\begin{frame}
  \frametitle{Анализ существующих решений}
  \begin{itemize}
    \item \textbf{Boost Graph Library (BGL)}
      \begin{itemize}
        \item (+) Широкий набор алгоритмов
        \item (-) Сложный API, высокий порог вхождения
        \item (-) Ограниченная поддержка многопоточности
      \end{itemize}
    \item \textbf{LEMON (Library for Efficient Modeling and Optimization in Networks)}
      \begin{itemize}
        \item (+) Хорошо подходит для задач комбинаторной оптимизации
        \item (-) Отсутствие нативной поддержки параллельных вычислений
      \end{itemize}
    \item \textbf{igraph с обёрткой для C++}
      \begin{itemize}
        \item (+) Функциональность для анализа социальных сетей
        \item (-) Зависимость от базовой C-библиотеки
        \item (-) Ограниченное использование современных возможностей C++
      \end{itemize}
  \end{itemize}
\end{frame}

% NEW. Функциональные требования
\section{Функциональные требования}
\begin{frame}
  \frametitle{Функциональные требования}
  \begin{itemize}
    \item \textbf{Представление графов} 
      \begin{itemize}
        \item Поддержка ориентированных и неориентированных графов
        \item Различные способы хранения 
        \item Поддержка взвешенных рёбер и метаданных вершин
      \end{itemize}
    \item \textbf{Алгоритмы}
      \begin{itemize}
        \item Базовые: BFS, DFS, топологическая сортировка
        \item Поиск кратчайших путей: Дейкстра, Беллман-Форд, A*
        \item Минимальные остовные деревья: Прим, Крускал
        \item Компоненты сильной связности, мосты, точки сочленения
      \end{itemize}
    \item \textbf{Многопоточность}
      \begin{itemize}
        \item Параллельные версии алгоритмов
        \item Настраиваемое количество потоков
        \item Потокобезопасность при параллельных операциях
      \end{itemize}
    \item \textbf{Экспорт и визуализация} в формат DOT
  \end{itemize}
\end{frame}

% 5. Архитектура библиотеки
\section{Архитектура библиотеки}
\begin{frame}
  \frametitle{Общая архитектура библиотеки}
  \begin{columns}
    \column{0.65\textwidth}
    \begin{itemize}
      \item \textbf{Модульная структура}
        \begin{itemize}
          \item \texttt{structure/} — представления графов
          \item \texttt{algorithms/} — классические и параллельные алгоритмы
          \item \texttt{visualization/} — инструменты визуализации
          \item \texttt{tests/} — инструменты тестирования
        \end{itemize}
      \item \textbf{Принципы ООП}
        \begin{itemize}
          \item Полиморфизм 
          \item Инкапсуляция 
          \item Наследование 
        \end{itemize}
      \item \textbf{Метапрограммирование}
        \begin{itemize}
          \item Поддержка произвольных типов для вершин и ребер через шаблоны, проверка корректности типов на этапе компиляции, и др.
        \end{itemize}
    \end{itemize}
    \column{0.35\textwidth}
    \centering
    \includegraphics[width=\columnwidth]{example_arch}
    \\{\footnotesize }
  \end{columns}
\end{frame}

% 6. Структуры представления графов
\section{Структуры представления графов}
\begin{frame}
  \frametitle{Структуры представления графов}
  \begin{columns}
    \column{0.6\textwidth}
    \begin{itemize}
      \item \textbf{Список смежности}
        \begin{itemize}
          \item Оптимален для разреженных графов
          \item $O(n+m)$ памяти
        \end{itemize}
      \item \textbf{Матрица смежности}
        \begin{itemize}
          \item Быстрая проверка наличия ребра за $O(1)$
          \item $O(n^2)$ памяти
        \end{itemize}
      \item \textbf{Список рёбер}
        \begin{itemize}
          \item Компактное хранение $O(m)$
          \item Медленный поиск ребра за $O(m)$
        \end{itemize}
      \item \textbf{CSR (Compressed Sparse Row)}
        \begin{itemize}
          \item Оптимизирован для параллельной обработки
          \item Эффективное хранение $O(n+m)$
        \end{itemize}
    \end{itemize}
    \column{0.4\textwidth}
    \centering
    \includegraphics[width=\columnwidth]{example_arch}
    \\{\footnotesize }
  \end{columns}
\end{frame}

% 7. Подход к многопоточным вычислениям
\section{Подход к многопоточным вычислениям}
\begin{frame}
  \frametitle{Подход к многопоточным вычислениям}
  \begin{itemize}
    \item \textbf{Стратегии распараллеливания}
      \begin{itemize}
        \item Разделение данных (data parallelism)
        \item Параллельное исследование графа из разных начальных точек
        \item Атомарные операции для безопасного обновления общих данных
      \end{itemize}
    \item \textbf{Обеспечение потокобезопасности}
      \begin{itemize}
        \item Мьютексы для защиты критических секций
        \item Атомарные операции для счетчиков
        \item Локальные буферы для промежуточных результатов
      \end{itemize}
    \item \textbf{Оптимизации}
      \begin{itemize}
        \item Предварительное вычисление метаданных для оптимизации
        \item Адаптивное определение оптимального числа потоков
        \item Балансировка нагрузки между потоками
      \end{itemize}
  \end{itemize}
\end{frame}

% NEWNEW. Выбор средств реализации
\section{Выбор средств реализации}
\begin{frame}
  \frametitle{Выбор средств реализации}
  \begin{itemize}
    \item \textbf{C++ (C++20)}
      \begin{itemize}
        \item Высокая производительность и контроль над памятью
        \item Шаблоны для обобщённого программирования
        \item Стандартная библиотека (STL) с контейнерами и алгоритмами
      \end{itemize}
    \item \textbf{Многопоточность}
      \begin{itemize}
        \item std::thread для создания потоков и управления ими
        \item Атомарные типы (std::atomic) для потокобезопасных операций
        \item Мьютексы (std::mutex) для синхронизации доступа к общим данным
      \end{itemize}
    \item \textbf{Тестирование}
      \begin{itemize}
        \item Google Test для модульных и интеграционных тестов
        \item Санитайзеры для выявления проблем (ThreadSanitizer)
      \end{itemize}
    \item \textbf{Визуализация:} GraphViz для построения графов по DOT-файлам
    \item \textbf{Сборка и документация:} CMake
  \end{itemize}
\end{frame}

% 8. Результаты тестирования производительности
\section{Результаты тестирования}
\begin{frame}
  \frametitle{Результаты тестирования производительности}
  \begin{columns}
    \column{0.48\textwidth}
    \textbf{BFS на большом графе}
    \begin{table}
    \centering
    \small  % Уменьшаем размер шрифта
    \begin{tabular}{|l|c|c|}
    \hline
    \textbf{Вариант} & \textbf{Время (с)} & \textbf{Уск.} \\
    \hline
    Послед. & 5.724 & 1.00x \\
    \hline
    Паралл. & 0.283 & 20.23x \\
    \hline
    \end{tabular}
    \end{table}

    \textbf{Параметры графа}
    \begin{itemize}
      \item 10M вершин
      \item 100M рёбер
      \item 8 потоков
    \end{itemize}

    \column{0.48\textwidth}
    \textbf{Алгоритм Дейкстры}
    \begin{table}
    \centering
    \small  % Уменьшаем размер шрифта
    \begin{tabular}{|l|c|c|c|}
    \hline
    \textbf{Размер графа} & \textbf{Посл.} & \textbf{Пар.} & \textbf{Уск.} \\
    \hline
    1M в., 10M р. & 9.48с & 6.76с & 1.40x \\
    \hline
    1M в., 100M р. & 70.69с & 8.20с & 8.62x \\
    \hline
    \end{tabular}
    \end{table}

    \textbf{Выводы}
    \begin{itemize}
      \item Значительный прирост при увеличении плотности графа
      \item CSR формат даёт лучшую локальность данных
    \end{itemize}
  \end{columns}
\end{frame}

% 9. Тестирование
\begin{frame}
  \frametitle{Тестирование разработанного решения}
  \begin{itemize}
    \item \textbf{Комплексная стратегия тестирования}
      \begin{itemize}
        \item Модульные тесты (Google Test)
        \item Интеграционные тесты
        \item Тесты производительности
      \end{itemize}
    \item \textbf{Тестирование корректности алгоритмов}
      \begin{itemize}
        \item Проверка на простых графах с известными результатами
        \item Тестирование граничных случаев
        \item Сравнение параллельных и последовательных реализаций
      \end{itemize}
    \item \textbf{Использование санитайзеров}
      \begin{itemize}
        \item ThreadSanitizer для обнаружения дата-рейсов и дэдлоков
        \item AddressSanitizer для выявления проблем с памятью
      \end{itemize}
    \item \textbf{Автоматизация тестирования}
      \begin{itemize}
        \item CMake для автоматической сборки проекта
        \item Автоматический запуск unit-тестов при сборке
        \item Интеграция с Google Test фреймворком
      \end{itemize}
  \end{itemize}
\end{frame}

% 10. Визуализация
\section{Визуализация}
\begin{frame}
  \frametitle{Визуализация графов}
  \begin{itemize}
    \item \textbf{Модуль визуализации} преобразует графы в формат DOT
      \begin{itemize}
        \item Поддерживает все реализованные структуры графов
        \item Генерирует текстовое описание на языке DOT
        \item Интегрируется с GraphViz для создания изображений
      \end{itemize}
    \item \textbf{Возможности визуализации}
      \begin{itemize}
        \item Настройка цветов вершин и рёбер
        \item Выделение определённых путей/подграфов
        \item Отображение весов рёбер и меток вершин
        \item Автоматическая компоновка графа
      \end{itemize}
    \item \textbf{Примеры визуализации}
      \begin{itemize}
        \item Графы дорожных сетей
        \item Минимальные остовные деревья
        \item Визуализация результатов алгоритмов поиска путей
      \end{itemize}
  \end{itemize}
\end{frame}

\begin{frame}{Пример визуализации}
  \begin{columns}
    \column{0.5\textwidth}
    \centering
    \includegraphics[width=\textwidth]{MST_road_network.png}
    \vspace{0.3em}
    \textbf{Исходный граф дорожной сети}
    
    \column{0.5\textwidth}
    \centering
    \includegraphics[width=\textwidth]{MST_road_network_mst.png}
    \vspace{0.3em}
    \textbf{Минимальное остовное дерево}
  \end{columns}
  
  \vspace{0.5em}
  \centerline{Граф российских дорог и его минимальное остовное дерево}
\end{frame}

% 11. Примеры применения
%\section{Примеры применения}
%\begin{frame}
  %\frametitle{Примеры практических применений}
  %\begin{itemize}
    %\item \textbf{Анализ социальных связей (гипотеза "шести рукопожатий")}
      %\begin{itemize}
        %\item Граф: 3+ миллиона вершин, 23+ миллиона рёбер
        %\item Многопоточный BFS (8 потоков) для анализа расстояний
        %\item Результат: 95.7\% пользователей доступны в пределах 6 шагов
        %\item Среднее расстояние: 5.13 шагов
      %\end{itemize}
    %\item \textbf{Многопоточный PageRank на веб-графе}
      %\begin{itemize}
        %\item Датасет: SNAP Google-Web
        %\item Значительное ускорение за счёт распределения нагрузки
        %\item Эффективное ранжирование узлов в крупном графе
      %\end{itemize}
  %\end{itemize}

  %Эксперименты демонстрируют возможность эффективной обработки реальных графов большого размера с использованием разработанных методов.
%\end{frame}



\begin{frame}{Пример применения: гипотеза шести рукопожатий}
  \begin{itemize}
    \item \textbf{Датасет:} социальный граф Facebook (3,097,165 вершин, 23,667,394 рёбер)
    \item \textbf{Эксперимент:} многопоточный BFS с 8 потоками от 100 случайных вершин
    \item \textbf{Ключевые результаты}
    \begin{itemize}
      \item 95.70\% пользователей соединены не более чем 6 рукопожатиями
      \item Среднее расстояние между пользователями: 5.13
      \item Загрузка и подготовка графа: 21.8 секунд
      \item Полный анализ: 81.7 секунд (с 8 потоками)
    \end{itemize}
    \item \textbf{Распределение расстояний}
    \begin{itemize}
      \item 4 рукопожатия: 20.33\%
      \item 5 рукопожатий: 44.59\%
      \item 6 рукопожатий: 28.63\%
    \end{itemize}
  \end{itemize}
  % Здесь будет диаграмма распределения расстояний
\end{frame}





%\begin{frame}{Пример применения: гипотеза шести рукопожатий}
  % Здесь будет диаграмма распределения расстояний
  %\includegraphics[width=0.63\textwidth]{SIX_distance_pie_chart_simple.png}
%\end{frame}





\begin{frame}{Пример применения: гипотеза шести рукопожатий}
  % Здесь будет диаграмма распределения расстояний
  \includegraphics[width=0.84\textwidth]{SIX_distance_bar_horizontal.png}
\end{frame}





\begin{frame}{Пример применения: PageRank для веб-графа Google}
  \begin{itemize}
    \item \textbf{Датасет:} веб-граф Google (875,713 вершин, 5,105,039 рёбер)
    \item \textbf{Эксперимент:} параллельный PageRank с разным числом потоков
    \item \textbf{Результаты ускорения}
    \begin{center}
      \begin{tabular}{|c|c|c|}
        \hline
        \textbf{Потоки} & \textbf{Время (с)} & \textbf{Ускорение} \\
        \hline
        1 & 5.68 & 1.00x \\
        2 & 4.51 & 1.26x \\
        4 & 3.51 & 1.62x \\
        8 & 3.01 & 1.89x \\
        16 & 2.61 & 2.18x \\
        \hline
      \end{tabular}
    \end{center}
    \item Анализ ТОП-100 вершин по рангу показывает ключевые страницы
    \item Реализация оптимизирована для эффективной работы с разреженными графами
  \end{itemize}
  % Здесь будет график ускорения
\end{frame}





\begin{frame}{Пример применения: PageRank для веб-графа Google}
  % Здесь будет диаграмма распределения расстояний
  \includegraphics[width=0.84\textwidth]{PR_pagerank_histogram.png}
\end{frame}





\begin{frame}{Пример применения: PageRank для веб-графа Google}
  % Здесь будет диаграмма распределения расстояний
  \includegraphics[width=0.99\textwidth]{PR_pagerank_clusters.png}
\end{frame}





% 12. Результаты и выводы
\section{Результаты и выводы}
\begin{frame}
  \frametitle{Результаты и выводы}
  \begin{itemize}
    \item \textbf{Результаты работы}
      \begin{itemize}
        \item Создана универсальная библиотека для работы с графами на C++
        \item Реализованы различные представления графов (список смежности, CSR и др.)
        \item Имплементированы классические и параллельные алгоритмы
        \item Достигнуто значительное ускорение (до 20x для BFS)
        \item Разработаны инструменты визуализации для анализа данных
      \end{itemize}
    \item \textbf{Преимущества решения}
      \begin{itemize}
        \item Современный C++ (использование новейших стандартов)
        \item Интуитивно понятное использование
        \item Встроенная поддержка многопоточности
        \item Гибкая и расширяемая архитектура
        \item Широкий спектр алгоритмов
      \end{itemize}
  \end{itemize}
\end{frame}

\begin{frame}
  \frametitle{Дальнейшие направления развития}
  \begin{itemize}
    \item \textbf{Расширение функционала}
      \begin{itemize}
        \item Добавление новых алгоритмов (потоковых, для сопоставления и др.)
        \item Реализация распределенных алгоритмов для вычисления на кластерах
        \item Интеграция с GPU-ускорением для избранных алгоритмов
      \end{itemize}
    \item \textbf{Улучшение производительности}
      \begin{itemize}
        \item Дополнительная оптимизация критических участков кода
        \item Адаптивное определение параметров параллелизации
        \item Исследование новых структур данных для специфических задач
      \end{itemize}
    \item \textbf{Дополнительные инструменты}
      \begin{itemize}
        \item Расширение возможностей визуализации
        \item Добавление аналитических инструментов для оценки характеристик графов
        \item Разработка готовых решений для типичных задач на графах
      \end{itemize}
  \end{itemize}
\end{frame}

% 13. Список источников
%\section{Список источников}
%\begin{frame}[allowframebreaks]
  %\frametitle{Список использованных источников}
  %\small
  %\begin{thebibliography}{99}
    %\bibitem{cormen-algorithms} \textit{Кормен Т., Лейзерсон Ч., Ривест Р., Штайн К.} Алгоритмы: построение и анализ, 3-е издание. — М.: Вильямс, 2013.
    %\bibitem{wilson-graph-theory} \textit{Wilson R.J.} Introduction to Graph Theory, 5th edition. — Pearson, 2010.
    %\bibitem{williams-concurrency} \textit{Williams A.} C++ Concurrency in Action. — Manning Publications, 2012.
    %\bibitem{snap-dataset} SNAP Datasets: Stanford Large Network Dataset Collection, \url{https://snap.stanford.edu/data/}
    %\bibitem{networkrepository} Network Repository, \url{http://networkrepository.com}
    %\bibitem{boost} Boost Graph Library (BGL), \url{https://www.boost.org/doc/libs/release/libs/graph/}
    %\bibitem{lemon} LEMON Graph Library, \url{https://lemon.cs.elte.hu}
    %\bibitem{igraph} igraph – The network analysis package, \url{https://igraph.org/}
  %\end{thebibliography}
%\end{frame}



%\begin{frame}{Выводы}
  %\begin{itemize}
    %\item \textbf{Разработана универсальная библиотека} для работы с графами на C++:
    %\begin{itemize}
      %\item Гибкая архитектура с поддержкой 4 способов хранения графов
      %\item Реализация классических и параллельных алгоритмов
      %\item Удобная визуализация через экспорт в формат DOT (GraphViz)
    %\end{itemize}
    %\item \textbf{Экспериментально доказана эффективность} параллельных реализаций:
    %\begin{itemize}
      %\item До 20x ускорение для BFS на плотных графах
      %\item До 8.62x ускорение для алгоритма Дейкстры на плотных графах
      %\item 2.18x ускорение для PageRank на веб-графе Google
    %\end{itemize}
    %\item \textbf{Практическая применимость} продемонстрирована на реальных задачах:
    %\begin{itemize}
      %\item Подтверждение гипотезы шести рукопожатий в социальных сетях
      %\item Ранжирование страниц в веб-графе
    %\end{itemize}
  %\end{itemize}
%\end{frame}

\begin{frame}{Доступ к исходному коду}
  \begin{itemize}
    \item \textbf{Проект доступен как open source на GitHub:} 
      \begin{itemize}
        \item \url{https://github.com/kunikrubika05/graph-library-cpp}
      \end{itemize}
    \item \textbf{В репозитории доступна подробная документация:}
      \begin{itemize}
        \item Инструкция по установке и сборке библиотеки
        \item Руководство по использованию с примерами
        \item Описание алгоритмов и структур данных
      \end{itemize}
  \end{itemize}
\end{frame}

% Заключение
%\begin{frame}
  %\centering
  %\Huge{Спасибо за внимание!}
%\end{frame}

\section{Результаты и выводы}
\begin{frame}
  \frametitle{Результаты и выводы}
  \begin{itemize}
    \item \textbf{Результаты работы}
      \begin{itemize}
        \item Создана универсальная библиотека для работы с графами на C++
        \item Реализованы различные представления графов (список смежности, CSR и др.)
        \item Имплементированы классические и параллельные алгоритмы
        \item Достигнуто значительное ускорение (до 20x для BFS)
        \item Разработаны инструменты визуализации для анализа данных
      \end{itemize}
    \item \textbf{Преимущества решения}
      \begin{itemize}
        \item Современный C++ (использование новейших стандартов)
        \item Интуитивно понятное использование
        \item Встроенная поддержка многопоточности
        \item Гибкая и расширяемая архитектура
        \item Широкий спектр алгоритмов
      \end{itemize}
  \end{itemize}
\end{frame}